% --- CORRECCION: Forzar interpretacion raw de la entrada ---
\UseRawInputEncoding
\documentclass[aspectratio=169, 10pt, xcolor=table]{beamer}

% --- Configuracion del Tema ---
\usetheme{Madrid}
\usecolortheme{dolphin}

% --- Paquetes necesarios ---
\usepackage[utf8]{inputenc}
\usepackage[spanish]{babel}
\usepackage{amsmath, amssymb}
\usepackage{booktabs}
\usepackage{graphicx}
\usepackage{listings}
\usepackage{tikz}
\usetikzlibrary{arrows.meta, positioning, shapes.geometric}
\usepackage{etoolbox}

% --- Ruta de Imagenes ---
\graphicspath{{./imagenes/}}

% --- Estilo para codigo Python ---
\definecolor{codegreen}{rgb}{0,0.6,0}
\definecolor{codegray}{rgb}{0.5,0.5,0.5}
\definecolor{codepurple}{rgb}{0.58,0,0.82}
\definecolor{backcolour}{rgb}{0.95,0.95,0.92}
\lstset{
    language=Python,
    backgroundcolor=\color{backcolour},
    commentstyle=\color{codegreen},
    keywordstyle=\color{blue},
    numberstyle=\tiny\color{codegray},
    stringstyle=\color{codepurple},
    basicstyle=\ttfamily\scriptsize,
    breaklines=true,
    frame=single,
    showstringspaces=false,
    literate={á}{{\'a}}1 {é}{{\'e}}1 {í}{{\'i}}1 {ó}{{\'o}}1 {ú}{{\'u}}1 {ñ}{{\~n}}1 {ü}{{\"u}}1
}

% --- Pie de pagina personalizado ---
\makeatletter
\setbeamertemplate{footline}{
  \leavevmode%
  \hbox{%
  \begin{beamercolorbox}[wd=.333333\paperwidth,ht=2.25ex,dp=1ex,center]{author in head/foot}%
    \usebeamerfont{author in head/foot}\insertshortauthor
  \end{beamercolorbox}%
  \begin{beamercolorbox}[wd=.333333\paperwidth,ht=2.25ex,dp=1ex,center]{title in head/foot}%
    \usebeamerfont{title in head/foot}Sesion 11
  \end{beamercolorbox}%
  \begin{beamercolorbox}[wd=.333333\paperwidth,ht=2.25ex,dp=1ex,right]{date in head/foot}%
    \usebeamerfont{date in head/foot}CALCULO Y ALGEBRA LINEAL \hfill \insertframenumber/\inserttotalframenumber \hspace*{0.3cm}
  \end{beamercolorbox}}%
  \vskip0pt%
}
\makeatother

% --- Datos de la portada ---
\title[SESION 11 CALCULO]{\textbf{Sesion 11}}
\subtitle{Determinantes, Inversas y Sistemas Lineales}
\author[Prof. C. Sanchez]{Carlos Cesar Sanchez Coronel}
\institute{\textbf{CALCULO Y ALGEBRA LINEAL PARA IA}}
\date{}

% --- Eliminar barra de navegacion ---
\setbeamertemplate{navigation symbols}{}

% --- Indice al inicio de cada seccion ---
\AtBeginSection[]{
  \begin{frame}
    \frametitle{Contenido}
    \tableofcontents[currentsection]
  \end{frame}
}

\begin{document}

% --- Portada ---
\begin{frame}
    \titlepage
\end{frame}

% --- Indice general ---
\begin{frame}{Contenido}
    \tableofcontents
\end{frame}

% ============================================================================
% SECCION 1: DETERMINANTE DE UNA MATRIZ
% ============================================================================
\section{Determinante de una Matriz}

\begin{frame}{Definicion e interpretacion geometrica}
    \begin{itemize}
        \item El determinante $\det(A)$ de una matriz cuadrada $A$ mide el cambio de volumen que induce la transformacion lineal asociada.
        \item $|\det(A)|$ es el factor por el cual se escala el area (2D) o volumen (3D).
        \item $\det(A) > 0$: preserva orientacion.
        \item $\det(A) < 0$: invierte orientacion (refleja).
        \item $\det(A) = 0$: la transformacion aplasta el espacio (perdida de dimension) $\Rightarrow$ matriz singular.
    \end{itemize}
    \centering
    \includegraphics[width=0.4\textwidth]{determinante_volumen.png} % si existe
\end{frame}

\begin{frame}{Calculo para matrices pequenas}
    \begin{block}{Matriz $2\times2$}
        \[
        \det\begin{pmatrix}a&b\\c&d\end{pmatrix} = ad - bc
        \]
    \end{block}
    \begin{block}{Matriz $3\times3$ (regla de Sarrus)}
        \begin{align*}
            \det\begin{pmatrix}
                a_{11} & a_{12} & a_{13} \\
                a_{21} & a_{22} & a_{23} \\
                a_{31} & a_{32} & a_{33}
            \end{pmatrix}
            &= a_{11}a_{22}a_{33} + a_{12}a_{23}a_{31} + a_{13}a_{21}a_{32} \\
            &\quad - a_{13}a_{22}a_{31} - a_{11}a_{23}a_{32} - a_{12}a_{21}a_{33}
        \end{align*}
    \end{block}
\end{frame}

\begin{frame}{Ejemplos de calculo}
    \begin{exampleblock}{Ejemplo 1}
        $\det\begin{pmatrix}3&5\\2&4\end{pmatrix} = 3\cdot4 - 5\cdot2 = 12-10 = 2$
    \end{exampleblock}
    \begin{exampleblock}{Ejemplo 2}
        $\det\begin{pmatrix}1&2&3\\0&1&4\\5&6&0\end{pmatrix} = 1\cdot1\cdot0 + 2\cdot4\cdot5 + 3\cdot0\cdot6 - (3\cdot1\cdot5 + 1\cdot4\cdot6 + 2\cdot0\cdot0) = 40 - 39 = 1$
    \end{exampleblock}
\end{frame}

\begin{frame}{Propiedades fundamentales}
    \begin{enumerate}
        \item $\det(I) = 1$.
        \item $\det(A^T) = \det(A)$.
        \item $\det(AB) = \det(A)\det(B)$.
        \item Si se intercambian dos filas, el determinante cambia de signo.
        \item Si una fila se multiplica por $k$, el determinante se multiplica por $k$.
        \item Si una fila es combinacion lineal de otras, $\det(A) = 0$.
        \item $\det(A^{-1}) = 1/\det(A)$ (si $A$ es invertible).
    \end{enumerate}
\end{frame}

% ============================================================================
% SECCION 2: MATRIZ INVERSA
% ============================================================================
\section{Matriz Inversa}

\begin{frame}{Definicion}
    \begin{block}{Matriz inversa}
        $A\in\mathbb{R}^{n\times n}$ es invertible si existe $A^{-1}$ tal que
        \[
        AA^{-1} = A^{-1}A = I_n
        \]
    \end{block}
    \begin{itemize}
        \item La inversa es unica.
        \item Existe si y solo si $\det(A) \neq 0$ (matriz no singular).
    \end{itemize}
\end{frame}

\begin{frame}{Condiciones equivalentes a invertibilidad}
    Son equivalentes:
    \begin{itemize}
        \item $\det(A) \neq 0$.
        \item $A$ tiene rango maximo ($=n$).
        \item Las columnas (y filas) de $A$ son linealmente independientes.
        \item El sistema $Ax=b$ tiene solucion unica para cualquier $b$.
        \item $A$ no tiene autovalor cero.
    \end{itemize}
\end{frame}

\begin{frame}{Calculo de la inversa}
    \begin{block}{Matriz $2\times2$}
        \[
        A = \begin{pmatrix}a&b\\c&d\end{pmatrix},\quad
        A^{-1} = \frac{1}{ad-bc}\begin{pmatrix}d&-b\\-c&a\end{pmatrix}
        \]
    \end{block}
    \begin{exampleblock}{Ejemplo}
        Para $A = \begin{pmatrix}4&7\\2&6\end{pmatrix}$, $\det = 24-14=10$,
        \[
        A^{-1} = \frac{1}{10}\begin{pmatrix}6&-7\\-2&4\end{pmatrix} = \begin{pmatrix}0.6&-0.7\\-0.2&0.4\end{pmatrix}
        \]
    \end{exampleblock}
\end{frame}

% ============================================================================
% SECCION 3: RANGO DE UNA MATRIZ Y NULIDAD
% ============================================================================
\section{Rango de una Matriz y Nulidad}

\begin{frame}{Definiciones}
    \begin{itemize}
        \item \textbf{Rango} de $A\in\mathbb{R}^{m\times n}$: dimension del espacio columna (o fila) = numero maximo de columnas (filas) linealmente independientes.
        \item \textbf{Nulidad}: dimension del espacio nulo $\{x: Ax=0\}$.
        \item \textbf{Teorema de la dimension}: $\operatorname{rango}(A) + \operatorname{nulidad}(A) = n$ (numero de columnas).
    \end{itemize}
\end{frame}

\begin{frame}{Interpretacion geometrica}
    \begin{itemize}
        \item El rango indica cuantas direcciones independientes puede generar la transformacion.
        \item Si $\operatorname{rango}(A) < n$, la transformacion aplasta el espacio en un subespacio de menor dimension (perdida de informacion).
        \item En redes neuronales, el rango de la matriz de pesos determina la capacidad expresiva de la capa.
    \end{itemize}
\end{frame}

\begin{frame}{Calculo del rango}
    \begin{itemize}
        \item Por eliminacion gaussiana (numero de pivotes no nulos).
        \item Mediante SVD (valores singulares mayores que una tolerancia).
    \end{itemize}
    \begin{exampleblock}{Ejemplo}
        $C = \begin{pmatrix}1&2&3\\2&4&6\\3&6&9\end{pmatrix}$. La segunda fila es el doble de la primera, la tercera el triple. Todas las filas son dependientes $\Rightarrow$ rango = 1.
    \end{exampleblock}
\end{frame}

% ============================================================================
% SECCION 4: APLICACIONES EN IA Y ROBOTICA
% ============================================================================
\section{Aplicaciones en IA, Robotica y Redes Neuronales}

\begin{frame}{Machine Learning: normal multivariante}
    La densidad de la normal multivariante:
    \[
    \mathcal{N}(\mathbf{x}|\boldsymbol{\mu},\boldsymbol{\Sigma}) = \frac{1}{(2\pi)^{d/2}|\boldsymbol{\Sigma}|^{1/2}} \exp\left(-\frac{1}{2}(\mathbf{x}-\boldsymbol{\mu})^\top \boldsymbol{\Sigma}^{-1}(\mathbf{x}-\boldsymbol{\mu})\right)
    \]
    \begin{itemize}
        \item $\boldsymbol{\Sigma}$: matriz de covarianza (definida positiva $\Rightarrow$ invertible, $\det>0$).
        \item $|\boldsymbol{\Sigma}|$ aparece en el factor de normalizacion.
        \item $\boldsymbol{\Sigma}^{-1}$ es la matriz de precision.
    \end{itemize}
\end{frame}

\begin{frame}{Robotica: determinante del Jacobiano}
    \begin{itemize}
        \item Jacobiano $J(q)$ relaciona velocidades articulares $\dot{q}$ con velocidad del efector $v$:
        \[
        v = J(q)\dot{q}
        \]
        \item Para un robot no redundante, $\det(J)=0$ indica una \textbf{singularidad}: se pierden grados de libertad (el efector no puede moverse en ciertas direcciones).
        \item Cerca de una singularidad, el determinante es pequeno y pequenas velocidades articulares producen grandes velocidades en el efector (inestabilidad).
    \end{itemize}
\end{frame}

\begin{frame}{Redes neuronales: rango de la matriz de pesos}
    \begin{itemize}
        \item Capa fully connected: $y = Wx$.
        \item El rango de $W$ determina la dimension del subespacio alcanzable.
        \item Rango completo = maxima capacidad expresiva.
        \item Rango bajo puede indicar redundancia o limitaciones en la representacion.
        \item Inicializaciones cuidadosas y regularizacion ayudan a mantener rango alto.
    \end{itemize}
\end{frame}

% ============================================================================
% SECCION 5: PYTHON CON NUMPY
% ============================================================================
\section{Python: Calculos con NumPy}

\begin{frame}[fragile]{Determinante, inversa y rango}
    \begin{lstlisting}
import numpy as np

A = np.array([[2,1],[1,2]])
print("Matriz A:\n", A)

det_A = np.linalg.det(A)
print("det(A) =", det_A)  # 3.0

A_inv = np.linalg.inv(A)
print("Inversa:\n", A_inv)
print("A @ A_inv:\n", A @ A_inv)  # identidad

rango_A = np.linalg.matrix_rank(A)
print("Rango de A =", rango_A)    # 2

# Matriz singular
B = np.array([[1,2],[2,4]])
det_B = np.linalg.det(B)           # 0
rango_B = np.linalg.matrix_rank(B) # 1
    \end{lstlisting}
\end{frame}

\begin{frame}[fragile]{Resolver sistemas lineales: inversa vs solve}
    \begin{lstlisting}
import time
n = 1000
A = np.random.randn(n,n)
b = np.random.randn(n)

# Metodo 1: usando inversa (ineficiente, inestable)
start = time.time()
x_inv = np.linalg.inv(A) @ b
t_inv = time.time() - start

# Metodo 2: usando solve (factorizacion LU, recomendado)
start = time.time()
x_solve = np.linalg.solve(A, b)
t_solve = time.time() - start

print(f"Tiempo inversa: {t_inv:.4f}s, solve: {t_solve:.4f}s")
    \end{lstlisting}
\end{frame}

\begin{frame}[fragile]{Visualizacion del determinante como factor de escala}
    \begin{lstlisting}
import matplotlib.pyplot as plt

x = np.linspace(-1,1,20)
y = np.linspace(-1,1,20)
X, Y = np.meshgrid(x,y)
puntos = np.vstack([X.ravel(), Y.ravel()])

theta = np.radians(30)
A = np.array([[2*np.cos(theta), -0.5*np.sin(theta)],
              [2*np.sin(theta),  0.5*np.cos(theta)]])
det_A = np.linalg.det(A)
puntos_trans = A @ puntos

plt.figure(figsize=(8,4))
plt.subplot(1,2,1)
plt.scatter(puntos[0,:], puntos[1,:], s=1)
plt.axis('equal'); plt.title('Original')
plt.subplot(1,2,2)
plt.scatter(puntos_trans[0,:], puntos_trans[1,:], s=1)
plt.axis('equal'); plt.title(f'Transformado (det = {det_A:.2f})')
plt.show()
    \end{lstlisting}
\end{frame}

\begin{frame}[fragile]{Rango y nulidad en red neuronal simulada}
    \begin{lstlisting}
W = np.random.randn(5,10)
print("Rango inicial:", np.linalg.matrix_rank(W))  # 5

# Hacer una columna dependiente
W[:,3] = 0.5*W[:,1] + 0.7*W[:,5]
print("Rango despues:", np.linalg.matrix_rank(W))  # 4

nulidad = W.shape[1] - np.linalg.matrix_rank(W)
print("Nulidad =", nulidad)
    \end{lstlisting}
\end{frame}

% ============================================================================
% SECCION 6: NOTACION EN PAPERS
% ============================================================================
\section{Notacion en Papers Cientificos}

\begin{frame}{Notacion comun}
    \begin{itemize}
        \item $\det(\boldsymbol{\Sigma})$: determinante de la matriz de covarianza.
        \item $\boldsymbol{\Sigma}^{-1}$: matriz de precision.
        \item $\det(\mathbf{J})$: determinante del Jacobiano (singularidades).
        \item $\operatorname{rank}(\mathbf{X})$: rango de la matriz de caracteristicas.
        \item En deep learning: $\operatorname{rank}(\mathbf{W}^{(l)})$ para estudiar capacidad expresiva.
    \end{itemize}
    \begin{exampleblock}{Ejemplo}
        Función de densidad normal multivariante:
        \[
        p(x) = \frac{1}{\sqrt{(2\pi)^d \det(\Sigma)}} \exp\left(-\frac{1}{2}(x-\mu)^\top \Sigma^{-1}(x-\mu)\right)
        \]
    \end{exampleblock}
\end{frame}

% ============================================================================
% SECCION 7: EJERCICIOS PROPUESTOS
% ============================================================================
\section{Ejercicios Propuestos}

\begin{frame}{Ejercicios}
    \begin{enumerate}
        \item Calcula el determinante, inversa (si existe) y rango de la matriz $\begin{pmatrix}0&1&4\\5&6&0\\1&2&3\end{pmatrix}$ usando NumPy y verifica.
        \item Demuestra que si $A$ es ortogonal ($A^\top A = I$), entonces $\det(A) = \pm 1$.
        \item Genera una matriz aleatoria $5\times5$ y otra multiplicándola por una matriz diagonal con un cero en la diagonal. Compara rangos y determinantes.
        \item Investiga como se calcula el rango usando SVD. Implementa una funcion que compute el rango con tolerancia.
        \item En un sistema de recomendacion, se factoriza $R \approx UV^\top$ con $U\in\mathbb{R}^{m\times k}$, $V\in\mathbb{R}^{n\times k}$. ¿Cual es el rango maximo de la aproximacion? ¿Por que se elige $k$ pequeño?
    \end{enumerate}
\end{frame}

% ============================================================================
% SECCION 8: CASO PRACTICO INTEGRADOR (JACOBIANO EN ROBOTICA)
% ============================================================================
\section{Caso Practico: Jacobiano en Robotica}

\begin{frame}{Problema}
    \begin{itemize}
        \item Un robot plano tiene Jacobiano $J = \begin{pmatrix}2&1\\1&2\end{pmatrix}$.
        \item Determinar si la configuracion es singular.
        \item Calcular la velocidad articular necesaria para lograr $v = (3,3)$.
        \item Interpretar que ocurre si $\det(J)$ fuese cercano a cero.
    \end{itemize}
\end{frame}

\begin{frame}{Solucion}
    \begin{block}{a) Determinante}
        $\det(J) = 2\cdot2 - 1\cdot1 = 3 \neq 0$ $\Rightarrow$ no es singular.
    \end{block}
    \begin{block}{b) Velocidad articular}
        $J^{-1} = \frac{1}{3}\begin{pmatrix}2&-1\\-1&2\end{pmatrix}$, $v=(3,3)$:
        \[
        \dot{q} = J^{-1}v = \frac{1}{3}\begin{pmatrix}2&-1\\-1&2\end{pmatrix}\begin{pmatrix}3\\3\end{pmatrix} = \frac{1}{3}\begin{pmatrix}3\\3\end{pmatrix} = \begin{pmatrix}1\\1\end{pmatrix}
        \]
    \end{block}
    \begin{block}{c) Si $\det(J)\approx 0$}
        Pequeñas velocidades articulares producirían grandes velocidades en el efector (inestabilidad). Se debe evitar trabajar cerca de singularidades.
    \end{block}
\end{frame}

% ============================================================================
% SECCION 9: RESUMEN
% ============================================================================
\section{Resumen}

\begin{frame}{Lo que hemos aprendido}
    \begin{exampleblock}{Conceptos clave}
        \begin{itemize}
            \item \textbf{Determinante}: mide cambio de volumen, condicion de invertibilidad.
            \item \textbf{Inversa}: permite deshacer transformaciones (¡no usar para resolver sistemas!).
            \item \textbf{Rango y nulidad}: dimension efectiva de una transformacion, capacidad expresiva.
            \item \textbf{Aplicaciones}: normal multivariante (ML), Jacobiano (robotica), rango en redes neuronales.
            \item \textbf{Python}: `np.linalg.det`, `inv`, `matrix_rank`, `solve`.
            \item \textbf{Notacion}: $\det(\Sigma)$, $\Sigma^{-1}$, $\operatorname{rank}(W)$.
        \end{itemize}
    \end{exampleblock}
\end{frame}

\begin{frame}{Conexion con futuros cursos}
    \begin{itemize}
        \item \textbf{Machine Learning}: analisis de covarianza, PCA (autovalores).
        \item \textbf{Robotica}: cinematica diferencial, control.
        \item \textbf{Deep Learning}: estudio de la capacidad de capas, regularizacion.
        \item \textbf{Optimizacion}: condiciones de segundo orden (hessiano definido positivo).
    \end{itemize}
\end{frame}

% --- Diapositiva final ---
\begin{frame}
    \centering
    \Huge \textbf{Gracias!}

\end{frame}

\end{document}